\section{Rating-Based User Feedback}
\label{sec:rating-based-user-feedback}

User preferences can be captured explicitly through numerical ratings, or 
implicitly by monitoring user behaviour 
\parencite[Chapter~3.1]{bobadilla2013recommender_systems_survey}. 
This feedback can be stored and used to improve recommendations 
in later interactions, so that the longer users engage 
with the system, the more the output can be refined to match 
their preferences 
\parencite[Chapter~1]{ricci2010recommender_systems_handbook}.

A practical problem with explicit ratings is that some items 
receive few ratings, making their average score unreliable. 
This is related to the sparse data problem, which frequently 
occurs when estimating counts from small amounts of data 
\parencite[Chapter~3.3]{murphy2012machine_learning_probabilistic_perspective}. 
Bayesian smoothing addresses this by pulling the estimated score 
towards a neutral prior when little data is available. 
The posterior is a compromise between what was previously 
believed and what the data indicates 
\parencite[Chapter~3.3]{murphy2012machine_learning_probabilistic_perspective}. 
As more data accumulates, the observed average comes to dominate 
the estimate, since the data has overwhelmed the prior 
\parencite[Chapter~3.3]{murphy2012machine_learning_probabilistic_perspective}.

The effectiveness of explicit feedback depends on having 
sufficient ratings. The cold start problem occurs when it is 
not possible to make reliable recommendations due to an 
initial lack of ratings 
\parencite[Chapter~3.2]{bobadilla2013recommender_systems_survey}. 
A particular variant is the new community problem, which refers 
to the difficulty of obtaining a sufficient amount of rating 
data when a system is first deployed 
\parencite[Chapter~3.2]{bobadilla2013recommender_systems_survey}. 
Under these conditions, enough ratings have to be collected 
before the system can understand user preferences and provide 
reliable recommendations 
\parencite[Chapter~3]{ricci2010recommender_systems_handbook}.   